% Category_圏論自主ゼミ
\documentclass[dvipdfmx]{jsreport}
\usepackage{amsmath, ascmac, latexsym, mathtools} %とりあえず入れとくパッケージ
% フォント関係
\usepackage{bm} % 小文字ベクトルのためのパッケージ
\usepackage{amssymb} % 集合の太文字パッケージ
\usepackage{mathrsfs} % 花文字を書くパッケージ
\usepackage{amsfonts}
\usepackage{fancybox} % 箱のとりあえずのデザイン(あんまり使わない)
\usepackage{amsthm} % 定理環境をカスタマイズ
\usepackage{tikz-cd} % 可換図式のパッケージ


\newcommand{\category}[1]{\mathcal{#1}}
\newcommand{\object}[1]{\mathrm{Ob}{(\category{#1})}}
\newcommand{\id}[1]{\mathrm{id}_{#1}}
\newcommand{\domain}[1]{\mathrm{dom}(#1)}
\newcommand{\codomain}[1]{\mathrm{cod}(#1)}

\begin{document}


\section{Introduction}
\subsection{圏論とは？}
    大学の講義で習ったような議論と同一の議論で集合，空間を扱う。

\subsection{前提}
    圏の定義として数学の公理系を仮定しないで行う(すなわちZFC公理系を仮定しない)。

    圏における対象とは集合全体や位相空間全体などを刺す。この対象の間には写像・連続写像・準同型写像などの射がある。それらは基本的に以下を満たす。

%\begin{enumerate_roman}
\begin{enumerate}
    \item 射(矢印)は合成可能
    \item 恒等写像が存在する
\end{enumerate}
%\end{enumerate_roman}
    これらの要素を抽出して圏を定義する。

%\begin{soudef}[圏]
    $\category{C}$が圏であるとは以下を満たすことである。
\begin{enumerate}
    \item $\category{C}$の対象とよばれるもの全体$\object{C}$が与えられている。
    \item $X, Y$を任意の2つの対象とする。この$X, Y$に対して$f$という$X$から$Y$への射または射の集合$\category{C}(X, Y)$が与えれている。
\end{enumerate}
%\end{soudef}
    このとき射$f$を$f\colon X\to Y$や$X\xrightarrow{f} Y$とかく。
    この射に対して以下3つを満たす。
\begin{enumerate}
    \item 合成可能性：
    $X, Y, Z$を対象とする。$f\colon X\to Y,\ g\colon Y\to Z$を射とする。このとき$X\xrightarrow{g\circ f} Z$をみたす。すなわち$g\circ f\in \category{C}(X, Z)$をみたす。

    \item 結合則：
    $X, Y, Z, W$を対象，$f\colon X\to Y,\ g\colon Y\to Z,\ h\colon Z\to W$とする。このとき$h\circ (g\circ f)= (h\circ g)\circ f$を満たす。

    \item 恒等写像の存在：
    $X, Y, Z$を対象とする。このときある射$\id{X}\in \category{C}$が存在して，$f\colon X\to Y,\ g\colon Y\to Z$に対して$f\circ \id{X}= f,\ \id{X}\circ g= g$を満たす。
\end{enumerate}

%\begin{souremark}
    射の集合は数学の集合の定義をみたすかわからない。
%\end{souremark}
%\begin{souremark}
    $X$が圏$\category{C}$の対象のとき，$X\in \object{C}$で表す。紛れがないときは$X\in \category{C}$とかき，$\id{X}$を1または$1_X$とかく。
%\end{souremark}

%\begin{souexample}
\begin{table}[h]
\begin{center}
\begin{tabular}{c|lll}
    記号& 圏の名前& 対象& 射\\ \hline
    Set& 集合& 集合& 写像\\
    set& 有限集合& 集合& 写像\\
    Ab& Abel群& Abel群& 準同型\\
    Top& 位相空間& 位相空間& 連続写像\\
\end{tabular}
\end{center}
\caption{考える対象と圏の名前}
\end{table}
%\end{souexample}

%\begin{soudef}[始域・終域]
    $\category{C}$を圏，$X, Y\in\object{C}$，$f\in \category{C}(X, Y)$とする。
    このとき$X$を$f$の始点または始域といい，$\domain{f}$とかく。同様に$Y$を$f$の終点または終域といい，$\codomain{f}$とかく。
%\end{soudef}

%\begin{souexample}
\begin{enumerate}
    \item 1つの対象とその対象の恒等写像からなる圏$\category{C}$
\begin{figure}

\end{figure}    

    \item 集合に対して対象全体を$S$とする。射の各元の恒等写像のみにすると
\begin{figure}

\end{figure}    
    のような圏が得られる。
\end{enumerate}
%\end{souexample}

\end{document}